\section{Lanczos and the leading direction}\label{kss:app:lanczos}

\fevc{} 0.6.0 finds the leading direction of Section~\ref{sec:q1} by a Lanczos iteration on the native route and
the matrix-free Mata route, where the prototype used power iteration; the exact Mata route uses a dense
eigendecomposition instead. Lanczos finds the extreme eigenvalues of a large symmetric operator $\mathcal M$ from products
$\mathcal Mv$ alone. It builds the Krylov subspace spanned by $r,\mathcal Mr,\mathcal M^2r,\dots$ from a start
vector $r$, orthonormalizes it with a three-term recurrence, and reads approximate eigenvalues (Ritz values) off
the small tridiagonal matrix that $\mathcal M$ becomes in that basis. The extreme eigenvalues converge first, at
a rate governed by the gap between $\lambda_1$ and $\lambda_2$ relative to the rest of the spectrum (roughly
$e^{-2k\sqrt{\mathrm{gap}}}$ after $k$ steps, against $(\lambda_2/\lambda_1)^k$ for power iteration); both ends of
the spectrum converge together, which matters for the covariance component, whose kernel has eigenvalues of both
signs. A thick restart keeps memory bounded.

In \fevc{} the problem is $Av_1=\lambda_1S_{xx}v_1$ of Section~\ref{sec:q1}, with $A$ the matrix of the component
and $S_{xx}=X'X$; its eigenvalues are those of the outcome-space kernel $B=XS_{xx}^{-1}AS_{xx}^{-1}X'$, and the
unit-space direction is $w_1=Xv_1$. The iteration runs in the $S_{xx}$ inner product, so each step is one linear
solve with $S_{xx}$ (the solver of the point estimate) plus a product with $A$; re-orthogonalization costs a
prediction and a transpose. The basis holds six vectors and a restart carries three Ritz vectors; the start vector
is a fixed function of the coefficient index (it never looks at the outcome) and each component starts from the
previous component's direction. The search stops when the residual of the leading pair is below the tolerance or
after 40 solves. When a direction dominates---the case that matters---the gap is large and a handful of steps
suffice; when no direction dominates, the top eigenvalues cluster, the search may stop at the cap, and the
direction carries a small share, so the rule picks the normal interval anyway. The second Ritz value is reported
because two large, close directions are the case that neither interval covers. KSS's code solves the same
generalized problem with \code{eigs}, which factors the explicit normal matrix and, for their non-symmetric form of
the problem, uses the Arnoldi-based Krylov--Schur method.
